% ============================================================
%  Section 3 -- Presentation de la solution
% ============================================================

\chapter{Présentation de la solution}

\section{Vue d'ensemble}

\bibopsname{} est organise comme un framework Python de benchmark. Le coeur du
projet est separe en quatre ensembles : l'agent IT, le moteur d'evaluation, les
pipelines de benchmark et la Racing Arena.

\begin{figure}[H]
  \centering
  \includegraphics[width=\textwidth]{architecture_overview.png}
  \caption{Vue d'ensemble de l'architecture \bibopsname{}.}
  \label{fig:architecture-overview}
\end{figure}

\begin{table}[H]
\centering
\begin{tabularx}{\textwidth}{lY}
\toprule
\textbf{Dossier} & \textbf{Responsabilite principale} \\
\midrule
\file{src/agent/} & Agent de support IT, boucle ReAct, outils, RAG, MCP. \\
\file{src/common/} & Configuration, clients LLM, helpers texte et extraction de tokens. \\
\file{src/bibops/evaluation/} & Juges, inspecteurs de securite, metriques composites, seuils. \\
\file{src/bibops/adapters/} & Adaptateurs vers agents IT, OpenAI-compatible, A2A et fact-checker. \\
\file{src/benchmark/} & Pipelines experimentaux et utilitaires de benchmark. \\
\file{src/racing/} & Hub FastAPI, moteur de course, clients LangGraph, equipes adversariales. \\
\file{scripts/} & Wrappers CLI et scripts de developpement. \\
\file{tests/} & Tests unitaires, fakes OpenAI et fixtures. \\
\bottomrule
\end{tabularx}
\caption{Organisation fonctionnelle du depot.}
\label{tab:repo-layout}
\end{table}

La base de code analysee contient 20\,326 lignes Python dans \file{src/},
\file{scripts/} et \file{tests/}. Cette taille reste compatible avec une
architecture lisible : les modules critiques sont courts, les tests isolent les
appels LLM, et les artefacts experimentaux sont stockes sous \file{data/}.

\section{Agent de support IT}

\subsection{Contrat de decision}

L'agent principal est implemente dans \file{src/agent/maestro.py}. La fonction
\code{lancer_agent()} execute une boucle ReAct bornee a cinq iterations par
defaut. A chaque tour, le LLM retourne un objet Pydantic \code{AgentDecision}
contenant trois champs :

\begin{itemize}
  \item \code{tool} : nom de l'outil a appeler ou \code{null} ;
  \item \code{argument} : argument normalise pour l'outil ;
  \item \code{final_answer} : reponse finale en francais lorsque plus aucun
  outil n'est necessaire.
\end{itemize}

Le choix d'un contrat structure est important. Il evite les parsings fragiles de
texte libre et rend les tests unitaires simples : les tests patchent
\code{_call_llm()} pour retourner directement une suite de decisions.

\subsection{Routage et memoire}

Avant le premier appel LLM, un routage lexical fournit un indice non bloquant.
Par exemple, les mots \code{vpn}, \code{cisco}, \code{outlook} ou \code{mail}
orientent vers \code{verifier_statut_serveur}, tandis que \code{bitlocker} ou
\code{procedure} orientent vers la documentation technique. Cet indice n'est
pas une contrainte : le modele conserve la decision finale.

La memoire courte, implemente dans \file{src/agent/memory.py}, conserve les
messages recents et les resultats outils. L'agent bloque aussi les appels
dupliques : si le modele redemande le meme outil avec le meme argument, le
systeme l'invite a produire une reponse finale.

\subsection{Outils disponibles}

Les outils sont declares dans \file{src/agent/tools.py}. Chaque outil possede
une politique de timeout, de retry et de longueur d'argument.

\begin{table}[H]
\centering
\begin{tabularx}{\textwidth}{lccY}
\toprule
\textbf{Outil} & \textbf{Timeout} & \textbf{Retries} & \textbf{Role} \\
\midrule
\code{verifier_statut_serveur} & 3 s & 0 & Interroge SQLite pour obtenir l'etat d'un service IT. \\
\code{chercher_dans_kb} & 5 s & 1 & Recherche lexicale dans la base de connaissances JSON. \\
\code{chercher_documentation_technique} & 8 s & 1 & Recherche RAG dans ChromaDB avec reranking hybride. \\
\bottomrule
\end{tabularx}
\caption{Politiques d'execution des outils de support.}
\label{tab:tools}
\end{table}

La recherche documentaire utilise des variantes de requete, une similarite
vectorielle et un score lexical. Les resultats sont filtres par
\code{RAG_DISTANCE_MAX = 1.2} ou par un recouvrement lexical minimal. Cette
hybridation est pragmatique : elle reduit les faux positifs d'une recherche
vectorielle pure lorsque les tickets sont courts.

\begin{figure}[H]
  \centering
  \includegraphics[width=\textwidth]{react_loop.png}
  \caption{Boucle ReAct de l'agent IT.}
  \label{fig:react-loop}
\end{figure}

\section{Moteur d'evaluation}

\subsection{Qualite}

Le module \file{src/bibops/evaluation/judges/} contient deux niveaux de juge.
\code{LLMJudge} est le primitif general : il note une reponse par rapport a un
critere. \code{LLMProfessor} specialise ce jugement pour les tickets IT, ajoute
du contexte RCA et persiste les resultats en SQLite lorsque necessaire.

Dans la campagne principale, le juge configure est \code{gpt-4o}. Le score
qualite est exprime sur 10. Il ne doit pas etre confondu avec le score composite
sur 100 : la qualite est une composante parmi d'autres, mais elle dispose aussi
d'un seuil bloquant.

\subsection{Securite}

La securite combine des detecteurs deterministes dans
\file{src/bibops/evaluation/checks.py}. Les familles couvertes sont :

\begin{itemize}
  \item donnees personnelles : email, telephone, IBAN, carte bancaire, adresse
  IPv4 ;
  \item secrets : cles OpenAI/Anthropic, GitHub PAT, AWS access key, bearer
  token, cle privee ;
  \item marqueurs de prompt injection en francais et en anglais ;
  \item refus explicites ;
  \item URLs suspectes ;
  \item marqueurs de toxicite simples.
\end{itemize}

Cette approche ne remplace pas une analyse semantique complete, mais elle a deux
avantages : elle est deterministe, rapide, testable sans reseau, et donne un
signal de gate exploitable dans un pipeline CI.

\subsection{Politique composite}

La classe \code{CompositePolicy} agrege les scores selon la formule :

\[
S = 0{,}40 Q + 0{,}35 Sec + 0{,}10 F + 0{,}10 L + 0{,}05 G,
\]

ou \(Q\) est la qualite normalisee, \(Sec\) la securite, \(F\) la normalisation
inverse du cout, \(L\) la normalisation inverse de la latence et \(G\) la
normalisation inverse de l'empreinte carbone. Le resultat est ensuite rapporte
sur 100.

\begin{figure}[H]
  \centering
  \includegraphics[width=0.72\textwidth]{composite_weights.png}
  \caption{Poids de la politique composite.}
  \label{fig:composite-weights}
\end{figure}

Les gates sont au moins aussi importants que la formule. Une architecture echoue
si le score qualite est inferieur a 7/10, si le score securite est inferieur a
6/10, si un blocage securite survient ou si les risques PII, injection,
non-refus ou toxicite depassent les seuils.

\section{Pipelines de benchmark}

La CLI \code{bibops} expose les commandes principales :
\code{bench}, \code{eval}, \code{dev} et \code{racing}. Les benchmarks
produisent des JSON sous \file{data/outputs/benchmark/} et des graphiques sous
\file{data/outputs/benchmark/charts/}. Les scripts peuvent aussi etre lances
directement avec \code{PYTHONPATH=.}.

Les pipelines principaux sont :

\begin{itemize}
  \item \code{bench compare-archs} pour comparer LLM unique et systeme
  multi-agents ;
  \item \code{bench ab-test} pour comparer deux modeles par jugement pairwise ;
  \item \code{bench position-bias} pour mesurer l'effet de l'ordre A/B ;
  \item \code{bench kaggle} pour rejouer une copie locale d'examen ;
  \item \code{dev mcp-server} et les benchmarks MCP pour exposer les outils ;
  \item \code{racing adversarial} pour lancer l'arene multi-agents.
\end{itemize}

\section{Racing Arena}

\subsection{Hub FastAPI et flux SSE}

La Racing Arena est une simulation distribuee. Le hub FastAPI, dans
\file{src/racing/hub/server.py}, expose notamment :

\begin{itemize}
  \item \code{GET /stream} : flux SSE de telemetrie commun a toutes les equipes ;
  \item \code{POST /decision/{team_id}} : reception des decisions ;
  \item \code{POST /ask_michelin} : RAG documentaire racing ;
  \item \code{GET /results} : historique complet des decisions ;
  \item des routes WeakProxy volontairement faibles pour les tests de securite.
\end{itemize}

Le moteur de course simule 15 tours dans la version observee, avec evolution de
la meteo, usure des pneus, carburant et possibilite de safety car. Le flux est
diffuse a tous les clients via des queues \code{asyncio}.

\begin{figure}[H]
  \centering
  \includegraphics[width=\textwidth]{racing_flux.png}
  \caption{Flux de donnees de la Racing Arena.}
  \label{fig:racing-flux}
\end{figure}

\subsection{Equipes et graphes LangGraph}

Chaque equipe est un processus separe. Le client lit le flux SSE, transforme la
telemetrie en etat \code{TeamState}, invoque un graphe LangGraph puis renvoie
une decision contrainte : \code{BOX BOX} ou \code{STAY OUT}. Les variantes
presentes dans le depot couvrent une equipe zero-shot, une equipe ReAct, une
equipe validee et Team Psi, l'attaquant.

Ce volet sert a tester des comportements qui n'apparaissent pas dans un simple
benchmark de tickets : reaction a un flux temps reel, effets d'une API non
authentifiee, attaques entre agents, fuite de strategie, et cout d'une
architecture multi-appels.

\section{Traces et reproductibilite}

L'agent IT peut sauvegarder des traces JSONL sous
\file{data/runtime/maestro/maestro_runs.jsonl}. Une trace contient l'identifiant
de run, les tours LLM, les appels outils, les durees, les previews de resultats,
la reponse finale, une reponse structuree et un niveau de confiance. Cette
instrumentation rend possible une analyse a posteriori : on peut relier un
score faible a une absence d'appel outil, a un timeout, a une source RAG
faiblement pertinente ou a un jugement qualite defavorable.

La reproductibilite repose aussi sur les tests. Les tests unitaires mockent les
LLM, notamment via \file{tests/_fakes/fake_openai.py} et des patches directs de
\code{_call_llm}. Cette strategie evite de rendre la qualite du test dependante
d'un service externe ou d'un modele instable.
